\chapter{KẾT LUẬN VÀ HƯỚNG PHÁT TRIỂN}
\label{ch:conclusion}

\section{Tổng kết Kết quả Đạt được}
\label{sec:conclusion_summary}

Bản đồ án tốt nghiệp với đề tài \textit{``Phương pháp tiếp cận Neuro-Symbolic trong việc sinh thủ tục kịch bản trinh thám có thể giải được sử dụng Hệ đa tác tử và Bộ kiểm định hình thức''} đã hoàn thiện mục tiêu thiết kế, cài đặt và đánh giá thực nghiệm đề ra. Đối chiếu với ba câu hỏi nghiên cứu xác lập tại Chương 1, đồ án thu được các kết quả cụ thể sau:

\begin{enumerate}[leftmargin=1.5cm]
    \item \textbf{Đối với RQ1 (Kiến trúc Phân rã Đa tác tử):} Tôi đã thiết kế và cài đặt hoàn chỉnh đường ống sinh kịch bản gồm 4 tác tử chuyên biệt chạy tuần tự (\texttt{OutlineAgent} $\to$ \texttt{CharacterAgent} $\to$ \texttt{ClueAgent} $\to$ \texttt{DeductionAgent}) kết hợp hàm tính cửa sổ thời gian tử vong (TOD) theo mô hình Henssge. Kiến trúc này giúp phân rã nhiệm vụ tạo sinh phức tạp thành từng lược đồ JSON hẹp, hạn chế hiện tượng suy giảm chú ý và tách bạch dữ liệu chân lý gốc khỏi quá trình chấm điểm phía người chơi.
    
    \item \textbf{Đối với RQ2 (Ràng buộc Ký hiệu và Tính Giải được):} Đồ án đã hình thức hóa các quy tắc trinh thám công bằng của Ronald Knox (1929), tích hợp phương trình làm lạnh hai hàm mũ Henssge (giải số bằng phương pháp chia đôi) cùng cấu trúc đồ thị hai phía vào bộ kiểm định \texttt{validate\_case\_diagnostics}. Qua 30 phiên thử nghiệm đối sánh trên 3 dòng mô hình LLM, hệ thống đạt 100\% tỷ lệ hợp lệ cấu trúc JSON và tỷ lệ vượt qua kiểm định logic từ 90\% đến 100\%, đồng thời chặn lọc chính xác ca lỗi tham chiếu thực thể phát sinh ở Lượt 5 của mô hình \texttt{gh/gpt-4o}.
    
    \item \textbf{Đối với RQ3 (Hiệu quả Sửa lỗi Cục bộ và Phản hồi Tất định):} Cơ chế sửa lỗi cục bộ định tuyến vi phạm về đúng tác tử chịu trách nhiệm với ngữ cảnh thu hẹp có độ phức tạp không gian token $\mathcal{O}(1)$ cho mỗi vi phạm (khoảng 300--500 tokens so với khoảng 2,500--3,500 tokens khi tái sinh toàn bộ chuỗi 4 tác tử). Bên cạnh đó, các thao tác kiểm định đồ thị, giải phương trình Henssge và chấm điểm kết án \texttt{evaluate\_player\_deduction} thực thi hoàn toàn trong bộ nhớ với thời gian dưới 5 mili-giây và 0 token, giữ độ trễ tạo sinh toàn trình ở phân vị P50 trong khoảng 10.73s--10.88s đối với \texttt{ollama/gemma4:31b-cloud} và \texttt{gh/gpt-4o-mini}.
    
    \item \textbf{Hoàn thiện Sản phẩm Đầu-cuối:} Toàn bộ kiến trúc đã được hiện thực hóa thành phần mềm chạy thực tế gồm máy chủ Python FastAPI (tích hợp xác thực HMAC-SHA-256, lưu trữ quan hệ SQLite WAL qua SQLAlchemy 2.0 và WebSocket chơi nhóm qua LAN) kết nối với Client Unity 6 URP 2D theo phong cách bảng điều tra trực quan.
\end{enumerate}

\section{Đóng góp Khoa học và Giá trị Thực tiễn}
\label{sec:conclusion_contributions}

\subsection{Giá trị về mặt Phương pháp luận}
Kết quả thực nghiệm của đồ án củng cố tính phù hợp của mô hình lai Neuro-Symbolic (Type 3 theo phân loại của Kautz) trong các ứng dụng tự sự tương tác đòi hỏi suy luận chặt chẽ. Thực nghiệm cho thấy việc giao phó hoàn toàn khâu kiểm chứng logic hoặc phán quyết kết án cho các LLM tự quy hồi vẫn tiềm ẩn rủi ro ảo giác tham chiếu thực thể; sự kết hợp giữa năng lực diễn đạt ngôn ngữ của LLM và màng lọc tất định bên ngoài là hướng tiếp cận thiết thực để nâng cao độ tin cậy của trò chơi giải đố.

\subsection{Giá trị Ứng dụng Thực tiễn}
Khung kiến trúc phân tách tầng tất định và tầng tạo sinh ngôn ngữ trong đồ án có thể mở rộng và tái sử dụng cho nhiều bài toán thực tiễn tương tự:
\begin{itemize}[leftmargin=1.5cm]
    \item Xây dựng hệ thống mô phỏng tình huống pháp lý hoặc diễn án giả định phục vụ đào tạo kỹ năng tranh tụng;
    \item Phát triển môi trường thực hành đọc hiểu hồ sơ, phân tích mốc thời gian y pháp và đối chiếu lời khai cho sinh viên ngành luật hoặc điều tra cơ bản;
    \item Ứng dụng vào các hệ thống hỗ trợ học tập thông minh, nơi LLM đóng vai trò hướng dẫn gợi ý bằng ngôn ngữ tự nhiên dựa trên kết quả chấm đúng/sai tất định của bộ giải toán học bên dưới.
\end{itemize}

\section{Các Hạn chế Hiện tại của Hệ thống}
\label{sec:conclusion_limitations}

Trên tinh thần đánh giá khách quan, hệ thống hiện tại vẫn còn bốn hạn chế kỹ thuật cần tiếp tục khắc phục:

\begin{enumerate}[leftmargin=1.5cm]
    \item \textbf{Chưa hỗ trợ thẩm vấn hội thoại tự do với NPC:} Tuân thủ phạm vi thiết kế dựa trên hồ sơ tĩnh để bảo đảm tính công bằng và ngăn chặn tấn công khai thác gợi ý trực tiếp, người chơi hiện chưa thể đặt câu hỏi ngôn ngữ tự nhiên tùy ý cho từng nghi phạm trong quá trình điều tra.
    
    \item \textbf{Kiến trúc lưu trữ đơn nút:} Dù đã triển khai cơ sở dữ liệu quan hệ SQLite (chế độ WAL với SQLAlchemy 2.0) hỗ trợ lưu trữ chuẩn ACID, hệ thống hiện vận hành ở mô hình máy chủ đơn lẻ và chưa hỗ trợ triển khai phân tán đa vùng khi mở rộng quy mô người chơi trực tuyến lớn.
    
    \item \textbf{Quy mô không gian kịch bản còn giới hạn:} Kịch bản sinh tự động hiện tập trung vào bối cảnh phòng kín với 3 nghi phạm và 3--4 vật chứng, chưa mở rộng sang các vụ án nhiều hiện trường liên thông hoặc mạng lưới đồng phạm phức tạp.
    
    \item \textbf{Giới hạn của prompt sửa lỗi đối với lỗi tham chiếu thực thể:} Như ghi nhận ở Lượt 5 của \texttt{gh/gpt-4o}, hàm \texttt{ClueAgent.repair\_clues} hiện ưu tiên bổ sung manh mối phá vỡ ngoại phạm và chưa chuyên biệt hóa cho tình huống sửa trường \texttt{linked\_suspect\_id} trỏ sai mã thực thể, khiến hồ sơ lỗi vẫn bị bộ kiểm định hậu kiểm giữ lại ở trạng thái không hợp lệ.
\end{enumerate}

\section{Hướng Nghiên cứu và Phát triển trong Tương lai}
\label{sec:conclusion_future_work}

Dựa trên nền tảng mã nguồn hiện có, các hướng nâng cấp tiếp theo được đề xuất gồm:

\subsection{Tích hợp Mô hình Ngôn ngữ Nhỏ Chạy Cục bộ}
Tiến hành tinh chỉnh các mô hình ngôn ngữ nhỏ mã nguồn mở (như dòng Llama-3 8B hoặc Qwen 2.5 7B) bằng kỹ thuật thích ứng bậc thấp LoRA (\textit{Low-Rank Adaptation}) trên tập hồ sơ vụ án đã vượt qua kiểm định Pydantic và Henssge. Giải pháp này hướng tới khả năng đóng gói mô hình chạy ngoại tuyến trên máy cá nhân nhằm giảm chi phí gọi API bên ngoài.

\subsection{Bổ sung Cơ chế Vá Cú pháp Tham chiếu và Đồ thị Tri thức Động}
Nâng cấp nhánh \texttt{repair\_clues} bằng thao tác sửa tất định trực tiếp trong mã nguồn đối với các lỗi \texttt{SCHEMA\_ERROR} đơn giản (hoặc ràng buộc danh sách mã hợp lệ trong prompt vá lỗi), đồng thời nghiên cứu sử dụng cơ sở dữ liệu đồ thị Neo4j với ngôn ngữ truy vấn Cypher để hỗ trợ kiểm tra tính giải được cho các kịch bản có nhiều nghi phạm và mắt xích chứng cứ bắc cầu dài hơn.

\subsection{Xây dựng Mô-đun Thẩm vấn Hội thoại Kiểm soát bằng FSM}
Phát triển tính năng đối thoại trực tiếp với NPC nghi phạm có kiểm soát trạng thái bằng FSM. Theo đó, câu trả lời của NPC do LLM sinh ra bị giới hạn trong tập dữ kiện cho phép tại từng trạng thái và được đối chiếu với trạng thái chân lý gốc trước khi hiển thị cho người chơi, vừa tăng tính tương tác vừa duy trì tính công bằng của trò chơi.
